%  RE-AUDITED 2026-07-06 (final pass, eleventh stop, light leg): zero
%  adjacency phrases -- position-portable at 11; trace-blind stitch
%  confirmed two-way and retargeted to the precise sec:born-rule label;
%  the quark-T3 routing line now names the frame (sec:frame) as the
%  defined term for the resolution classes; header umlaut normalised,
%  file pure ASCII; no other prose changes -- the deep audit stands.
% =====================================================================
%  Section: Quantum Numbers   (Plan Stage 2)
%  Paper 1, "It from Bit via G\"odel". Drafted 2026-06-08.
%  Follows \S sec:register; precedes \S sec:generations (the particle zoo).
%
%  \input fragment. biblatex/\autocite, amsthm `proposition'.
%  No new references (krasnov2021so9, szangolies2025standardmodel already
%  in the library; GHZ/W are framework-internal, Appendices B--C).
%
%  Cross-reference labels (placeholders -- re-point before compiling):
%    sec:register              -- the three-qubit register (Stage 1)
%    sec:entanglement-classes  -- GHZ/W theorems
%    sec:canonical-orbits      -- canonical orbit = particle (postulate)
%    prop:ghz, prop:w-doublet  -- those theorems (rescoped 2026-06-11)
%    prop:gsm                  -- G_SM as the Spin(9) stabiliser
%    app:ghz-proof, app:w-proof-- the proofs
%    sec:generations           -- the particle zoo / three generations (forward)
%    sec:boson-masses          -- boson masses, where sin^2 theta_W is applied (forward)
%    sec:interactions          -- the quantum-circuit branch (forward)
%
%  REVISED 2026-06-11 after review: canonical-orbit scope; corrected
%  doublet language; tracelessness flagged as the load-bearing assumption
%  (per hypercharge_derivation.md 7.2); anomaly claim made precise
%  (doublet-sector condition automatic, higher anomalies flagged);
%  right-handed singlet hypercharges added (Y = 2Q via the zigzag);
%  garbled SU(5) sentence fixed; matching scale standardised to 3.6 TeV.
%
%  OPEN-PROBLEM CONVENTION (paper-wide, introduced 2026-06-11): open
%  problems are surfaced where encountered as numbered environments with
%  op: labels, and collected in a closing "Open Problems" list that
%  references them. Preamble (option A, recommended -- automatic list):
%      \usepackage{thmtools}                     % after amsthm
%      \declaretheorem[style=definition,name=Open Problem]{openproblem}
%  and at the end of the paper:
%      \renewcommand{\listtheoremname}{Open Problems}
%      \listoftheorems[ignoreall,show={openproblem}]
%  (option B, manual: \theoremstyle{definition}
%  \newtheorem{openproblem}{Open Problem} and a hand-written closing
%  section using \ref{op:...}.)
%  Labels in this section: prop:hypercharge, prop:trace-blind (added
%  2026-07-04), op:traceless (shrunk 2026-07-04), op:anomalies,
%  sec:rh-singlets.
%
%  Hypercharge convention: Q = T_3 + Y/2, with Y(lepton C-singlet) = -1 and
%  Y(quark C^3-triplet) = +1/3, matching \S sec:entanglement-classes.
% =====================================================================

\section{Quantum Numbers}
\label{sec:charges}

The register of \S\ref{sec:register} carries three qubits, one per gauge sector, and the theorems
of \S\ref{sec:entanglement-classes} have already told us which canonical entanglement orbit is a
quark and which a lepton --- the identification of \S\ref{sec:canonical-orbits}. What remains is to
read the actual quantum numbers --- colour, weak isospin,
hypercharge, electric charge --- off a state, and to see that they come out quantised at the
Standard Model values rather than assigned to fit them. This section does that, recovers the
Gell-Mann--Nishijima relation as an identity of the register, and derives the weak mixing angle.
Throughout, a charge is to be understood operationally: it is what a gate acting on the relevant
qubit reads or changes, a framing we use lightly here and lean on when interactions are built as
circuits in \S\ref{sec:interactions}.

\subsection{Colour and weak isospin from the entanglement classes}

The two non-abelian charges are already in hand. By Proposition~\ref{prop:ghz} a quark channel ---
a canonical GHZ orbit, under the identification of \S\ref{sec:canonical-orbits} --- has its Hopf
image on the unit sphere of the colour-triplet sector $\mathbb{C}^3$, so its $C$-qubit carries the
$\mathbf{3}$ of $SU(3)$: this is colour. Two levels of one structure should not be conflated
here: the GHZ class that makes the channel a quark is the register's own --- the entanglement among
its $A$-, $B$- and $C$-qubits, which is what Proposition~\ref{prop:ghz} classifies --- while within
a hadron the quark's $C$-qubit is additionally GHZ-networked with its partners' $C$-qubits, the
closed colour structure of \S\ref{sec:confinement}. The species is register-level; the binding is
inter-channel; both are GHZ, one scale apart. A lepton channel --- a canonical W orbit --- is a colour singlet:
its $C$-qubit carries the trivial representation. Operationally, colour is what a gate on the
$C$-qubit reads, and the confinement of Proposition~\ref{prop:ghz} is, at register level, the
statement that no local gate moves a quark channel off the colour locus: no gate can isolate a
single colour. (The dynamical content of confinement belongs to Paper~2; only this structural form
is used here.)

Weak isospin is the content of Proposition~\ref{prop:w-doublet}: the single-qubit reduced state of
a W-class channel has the non-degenerate spectrum $(\tfrac23,\tfrac13)$, which singles out a
distinguished $T_3$ axis and an ordered pair of eigenstates --- the two doublet members --- so its
$B$-qubit carries $T_3 = \pm\tfrac12$ along that axis: for the leptons, the neutrino and the
charged lepton. The quark doublet carries the same $T_3 = \pm\tfrac12$ structure, supplied not by
the W-marginal spectrum but by the resolution classes of \S\ref{sec:generations} --- the frame of \S\ref{sec:frame}. A
weak-isospin
singlet has $T_3 = 0$. The raising and lowering operators
about the distinguished axis are exactly what a $W^\pm$ gate on the $B$-qubit applies, taking one
doublet member to the other; that the doublet is left-handed --- that only one chirality sits in it
--- is the handedness foreshadowed in \S\ref{sec:entanglement-classes} and made precise with the
generation structure in \S\ref{sec:generations}.

\subsection{Hypercharge from the trace condition}

Hypercharge is the one abelian charge. The $U(1)_Y$ generator must commute with the colour $SU(3)$,
and since $SU(3)$ acts irreducibly on $\mathbb{C}^3$ and trivially on $\mathbb{C}$, Schur's lemma
forces it to be diagonal on the split: a phase rotation acting on
$\mathbb{O} = \mathbb{C}\oplus\mathbb{C}^3$ with constant eigenvalues $y_{\mathbb{C}}$ on the
singlet (the lepton direction) and $y_{\mathbb{C}^3}$ on the colour triplet (the quark directions).
One further condition fixes the eigenvalues up to scale: the generator must be \emph{traceless}
over $\mathbb{O}$. Since $\mathbb{C}$ is one complex dimension and $\mathbb{C}^3$ is three,
\begin{equation}
  1\cdot y_{\mathbb{C}} + 3\cdot y_{\mathbb{C}^3} = 0
  \quad\Longrightarrow\quad
  y_{\mathbb{C}} = -3\,y_{\mathbb{C}^3}.
  \label{eq:trace}
\end{equation}

\begin{proposition}[Hypercharge spectrum]\label{prop:hypercharge}
Tracelessness over $\mathbb{O} = \mathbb{C}\oplus\mathbb{C}^3$ fixes the hypercharge eigenvalues on
the lepton and quark sectors to the ratio $y_{\mathbb{C}} : y_{\mathbb{C}^3} = -3 : 1$. Fixing the
sign and overall normalisation by the electron's charge gives $(y_{\mathbb{C}}, y_{\mathbb{C}^3}) =
(-1, +\tfrac13)$.
\end{proposition}

The status of the trace condition can now be stated more strongly than when this section was first
assembled, because half of the promotion has since become a theorem. A trace part would be a phase
rotation proportional to the identity on $\mathbb{O}$ --- an overall phase of the whole register.
That direction is provably invisible to records: the computation of
Proposition~\ref{prop:f1-dictionary} is dimension-independent and extends verbatim from
$\mathbb{C}^{2}$ to the register (the Born-rule construction of \S\ref{sec:born-rule} rests
on the same fact).

\begin{proposition}[The trace direction is record-blind]\label{prop:trace-blind}
Let the phase torus act on the register $\mathbb{O} = \mathbb{C}\oplus\mathbb{C}^{3}$, and let
the trace direction be its diagonal $U(1)$. Then every classical structure on the register
transports under the trace direction by a global scalar; every classical record any structure
induces is independent of it; and consequently every record-level holonomy --- every Wilson-loop
action on records --- assigned to the trace direction is the identity. An executable verification
of the finite checks, including the contrast that a traceless generator does act on records,
accompanies the code release.
\end{proposition}

\begin{proof}
The computation of Proposition~\ref{prop:f1-dictionary} in $n$ dimensions: for a classical
structure $\delta$ with copyables $\{v_{i}\}$, transport by $e^{i\alpha}\,\mathrm{id}$ gives
$\delta \mapsto e^{i\alpha}\delta$ with copyables $e^{i\alpha}v_{i}$ --- a global scalar ---
and the induced record channel
$\delta^{\dagger}(\rho\otimes I)\delta = \sum_{i}\langle v_{i}|\rho|v_{i}\rangle\,
v_{i}v_{i}^{\dagger}$ depends only on the projectors, hence not on $\alpha$. A holonomy whose
action on every record is trivial is the identity at record level.
\end{proof}

The second half of the promotion is then a one-paragraph conditional. Any gauge kinetic term that
is a functional of record-observable holonomy is, by Proposition~\ref{prop:trace-blind},
independent of the trace component of the connection: the trace mode has no field strength such a
functional can register, acquires no kinetic term, and does not propagate. For every kinetic term
of that kind, the gauge-observable $U(1)$ is the traceless part of the phase rotation on
$\mathbb{O}$. What remains open is exactly one construction, and the problem shrinks to it.

\begin{openproblem}[Tracelessness of $U(1)_Y$]\label{op:traceless}
Construct the gauge kinetic term from the Hopf holonomy explicitly, and establish that it is a
functional of record-observable holonomy, so that --- with Proposition~\ref{prop:trace-blind} and
the conditional argument above --- the gauge-observable $U(1)$ is the traceless part of the phase
rotation on $\mathbb{O}$ unconditionally. The record-blindness half is discharged; what is open is
the construction itself, which ties this problem to the loop-derived action of the dynamical layer.
\end{openproblem}

Granted it, there is no further freedom: the relative hypercharge of leptons and quarks is not an
input chosen to make the theory consistent but a consequence of one phase rotation being traceless
on the computability split.

The trace condition also does anomaly work, and here precision matters. With the real
multiplicities of the split, $2\,y_{\mathbb{C}} + 6\,y_{\mathbb{C}^3} = 0$ is exactly the vanishing
of the summed hypercharge of the left-handed doublet sector --- the mixed $SU(2)^2$--$U(1)_Y$
anomaly condition --- so that cancellation is automatic rather than imposed: the same condition
that fixes the eigenvalues is the condition that cancels it. The remaining conditions --- the cubic
$U(1)_Y^3$, the mixed $SU(3)^2$--$U(1)_Y$, and the mixed gauge--gravitational condition --- involve
the right-handed singlet hypercharges below; they are satisfied by the resulting assignments, as in
the Standard Model, but their structural status is open. One clarification is owed on the third,
whose name can mislead: despite the word, it is the least gravitational of the three. Its content
is the bare trace $\sum Y = 0$ over a complete generation; gravity enters only as the historical
origin of the requirement --- that the hypercharge current survive coupling to an arbitrary
background --- and since the background is, in this framework, itself emergent
(\S\ref{sec:space-of-computations}), the requirement is not optional here. It is also the nearest
of the three to discharged: given Proposition~\ref{prop:trace-blind} --- the gauged generator is
traceless --- and the completeness of the generation multiplet (\S\ref{sec:generations},
\S\ref{sec:zoo-singlets}), the spectrum sum is the trace of a traceless generator over a complete
multiplet, and the assignments above do sum to zero sector by sector: the doublets by the
$SU(2)^{2}$ condition already shown automatic, and the singlets on their own account.

\begin{openproblem}[Higher anomalies from traces]\label{op:anomalies}
Derive the cubic $U(1)_Y^3$, the mixed $SU(3)^2$--$U(1)_Y$ and the mixed gauge--gravitational
cancellations structurally, each as a trace condition on the appropriate sector of the register,
rather than as a check on the resulting assignments. By the paragraph above, the third reduces to
establishing that a generation is a complete multiplet of the traceless generator --- the linear
trace; the cubic and colour-weighted traces are the genuinely open pair.
\end{openproblem}

The right-handed singlets complete the table.\label{sec:rh-singlets} They have $T_3 = 0$, and the
chirality flip that produces them preserves electric charge (the \emph{Zitterbewegung} zigzag of Penrose, an
approach he himself calls ``slightly non-standard'' \autocite[628--632]{Penrose1994}, adopted here
and made circuit-level at the Higgs vertex of \S\ref{sec:interactions}), so their
hypercharges are fixed by $Y = 2Q$: $-2$ for the charged leptons, $+\tfrac43$ and $-\tfrac23$ for
the up- and down-type quarks, and $0$ for the sterile class of \S\ref{sec:generations}. This is the
full Standard Model assignment, with the zigzag mechanism flagged as adopted rather than
re-derived. As spectrum content --- channels with no frustrated pair for the weak force to act on
--- the singlets are enumerated with the rest of the one-generation table in
\S\ref{sec:zoo-singlets}.

\subsection{The Gell-Mann--Nishijima relation as a register identity}

Combining the $B$-qubit eigenvalue $T_3$ with the $A$-qubit eigenvalue $Y$ through
\begin{equation}
  Q = T_3 + \tfrac12 Y
\end{equation}
reproduces the electric charges of an entire generation, with the hypercharges of
Proposition~\ref{prop:hypercharge}:

\begin{center}
\begin{tabular}{lcccc}
\hline
state & sector & $T_3$ & $Y$ & $Q = T_3 + Y/2$ \\
\hline
neutrino $\nu$        & $\mathbb{C}$   & $+\tfrac12$ & $-1$          & $0$ \\
charged lepton $e$    & $\mathbb{C}$   & $-\tfrac12$ & $-1$          & $-1$ \\
up-type quark $u$     & $\mathbb{C}^3$ & $+\tfrac12$ & $+\tfrac13$   & $+\tfrac23$ \\
down-type quark $d$   & $\mathbb{C}^3$ & $-\tfrac12$ & $+\tfrac13$   & $-\tfrac13$ \\
\hline
\end{tabular}
\end{center}

The relation $Q = T_3 + \tfrac12 Y$ is not imposed: $T_3$ is the $SU(2)$ label of
Proposition~\ref{prop:w-doublet}, $Y$ the phase eigenvalue of Proposition~\ref{prop:hypercharge},
and $Q$ is the combination of the $A$- and $B$-qubit charges that the electromagnetic gate reads
(the factor $\tfrac12$ is a normalisation convention, absorbable into the scale of $Y$).
The Gell-Mann--Nishijima formula is thus an identity of the register.

\subsection{Charge quantisation and the \texorpdfstring{$\mathbb{Z}_6$}{Z6} quotient}

The fractional values --- third-integer charges for quarks, integer for leptons --- are quantised,
and the framework says why. The gauge group is not the bare product but the quotient
$G_{\mathrm{SM}} = \bigl(SU(3)\times SU(2)\times U(1)\bigr)/\mathbb{Z}_6$ established in
\S\ref{sec:entanglement-classes} as the $\mathrm{Spin}(9)$ stabiliser
(Proposition~\ref{prop:gsm}) \autocite{krasnov2021so9}. The $\mathbb{Z}_6$ identifies a combination
of the $\mathbb{Z}_3$ centre of colour, the $\mathbb{Z}_2$ centre of weak isospin, and a $U(1)_Y$
phase that acts trivially on every physical channel. For that element to act trivially, the
hypercharge of a channel must be locked to its colour triality and its weak isospin: a colour
triplet must carry hypercharge in units of $\tfrac13$, modulo integers, relative to a colour
singlet, which is exactly the offset between $y_{\mathbb{C}^3} = +\tfrac13$ and
$y_{\mathbb{C}} = -1$. The trace
condition~\eqref{eq:trace} fixes the ratio; the $\mathbb{Z}_6$ quotient forces it to be realised in
commensurate, quantised units. Electric charge is then quantised in thirds, with leptons landing on
integers because they are colour singlets --- the same structural fact, read through $Q = T_3 +
\tfrac12 Y$.

\subsection{The weak mixing angle}

The mixing of $U(1)_Y$ with the diagonal $SU(2)_L$ generator is set by the ratio of the two gauge
couplings,
\begin{equation}
  \sin^2\theta_W = \frac{g'^2}{g^2 + g'^2}.
\end{equation}
At tree level the framework fixes this ratio geometrically. The weak $SU(2)$ lives on the
three-dimensional fibre $S^3$ and the hypercharge $U(1)$ on the one-dimensional fibre $S^1$
(\S\ref{sec:register}); with no feature distinguishing one Lie-algebra direction from another, the
natural metric on the combined electroweak fibre weights every direction equally, so the
coupling-squared of each factor scales with its fibre dimension, $g^2 \propto \dim S^3 = 3$ and
$g'^2 \propto \dim S^1 = 1$. Hence
\begin{equation}
  \sin^2\theta_W = \frac{1}{1+3} = \tfrac14
  \qquad\text{(bare, tree level).}
\end{equation}
Equivalently: of the four electroweak generators, three engage the weak fibre and one survives
unbroken --- massless --- as the photon, so the photon occupies a fraction $\tfrac14$ of the
electroweak algebra. (This massless direction is the one \S\ref{sec:masses} calls \emph{debt-free},
once mass has been identified there with octonionic associator-debt; at this point we use only its
masslessness.)

This is a bare value, to be run to the electroweak scale by the ordinary Standard Model
renormalisation group, and its origin should not be confused with grand unification: $SU(5)$
predicts a bare $\tfrac38$ at its unification scale, whereas here there is no $SU(5)$ --- the
unified group does not arise from the division-algebra ladder, consistent with the non-observation
of proton decay --- and the bare value is $\tfrac14$ from the fibre dimensions. Running $\tfrac14$
down with the measured couplings reaches the observed $\sin^2\theta_W \approx 0.231$ at $M_Z$, the
bare value being attained at a matching scale of roughly $3.6\ \mathrm{TeV}$ --- about fifteen
times the Higgs vacuum value, and just above the scale at which the quaternionic structure becomes
physically operative.

Two honesty notes. The derivation rests on the structural assumption that the tree-level couplings
scale with the fibre dimension; we flag it as such, while noting that the same value $\tfrac14$ is
reached independently by counting the single massless direction, which is reassuring. And the
identification of \emph{which} combination is the photon --- the direction blind to the generation
structure as well as VEV-invariant --- is sharpened only once that structure is in place, in
\S\ref{sec:generations}; here we have used only that one of the four directions survives unbroken.
We derive $\sin^2\theta_W = \tfrac14$ at this point because it is fixed by the hypercharge embedding,
and apply it to the boson masses in \S\ref{sec:boson-masses}.

\subsection{Charges as operations}

For the circuit picture of \S\ref{sec:interactions} it is worth restating the quantum numbers
operationally, since that is the form interactions consume. Colour is what a gate on the $C$-qubit
reads; weak isospin is what a $W^\pm$ gate raises or lowers on the $B$-qubit; hypercharge is the
phase a gate applies to the $A$-qubit; and electric charge $Q = T_3 + \tfrac12 Y$ is what the
electromagnetic gate reads, a combination that mixes the $A$ (hypercharge) and $B$ (isospin)
channels --- which is why the photon is the surviving $A$--$B$ combination rather than a pure
$A$-qubit phase. Each gauge boson thus acts through a definite part of the register --- the
$W^\pm$ on the $B$-qubit, the gluons on $C$, the photon and the $Z$ on the $A$--$B$ pair jointly
--- and a channel's interactions are fixed by which of its three qubits carry non-trivial state.
This is the hook the interaction circuits build on; we develop it there.

With colour, weak isospin, hypercharge and electric charge all read off the register, and all
quantised at their measured values, a channel's full Standard Model quantum-number assignment is
now legible from its entanglement structure alone. The next step is to enumerate the channels ---
the fermion content of one generation, and then why there are three --- which is the business of
\S\ref{sec:generations}.
